\chapter{Three-Robot Cluster Kinematics}
\label{app:kin3}

% Ported from: cwaight_systems_paper.tex
\section{Forward Kinematics}

The forward kinematics transforms robot positions $(x_1, y_1)$, $(x_2, y_2)$, $(x_3, y_3)$ to the SAS parameters and centroid coordinates. The centroid position is:
\begin{equation}\label{eq:appB:20}
x_c = \frac{x_1 + x_2 + x_3}{3}, \quad y_c = \frac{y_1 + y_2 + y_3}{3}
\end{equation}
The side lengths, with $\mathbf{p}_i = (x_i, y_i)$ the position of robot $i$, are:
\begin{equation}\label{eq:appB:21}
p = \lVert \mathbf{p}_2 - \mathbf{p}_1 \rVert, \quad q = \lVert \mathbf{p}_3 - \mathbf{p}_2 \rVert
\end{equation}
The interior angle at robot 2 is:
\begin{equation}\label{eq:appB:23}
\beta = \arccos\left(\frac{(\mathbf{p}_1 - \mathbf{p}_2) \cdot (\mathbf{p}_3 - \mathbf{p}_2)}{pq}\right)
\end{equation}

The orientation angle is:
\begin{equation}\label{eq:appB:24}
\theta_c = \text{atan2}(y_1 - y_2, x_1 - x_2)
\end{equation}

\section{Inverse Kinematics}

Given the SAS parameters $(p, \beta, q)$ and centroid parameters $(x_c, y_c, \theta_c)$, the robot positions are recovered by constructing the triangle in a local frame and then applying rotation and translation.

In the local frame with robot 2 at the origin, the triangle is constructed as:
\begin{equation}\label{eq:appB:25}
\begin{aligned}
(x_1, y_1)^{\text{local}} &= (p,\, 0), \qquad (x_2, y_2)^{\text{local}} = (0,\, 0), \\
(x_3, y_3)^{\text{local}} &= (q\cos\beta,\, q\sin\beta)
\end{aligned}
\end{equation}

This construction assumes robot 3 lies counterclockwise from the robot 1--robot 2 edge; since $\beta$ from (\ref{eq:appB:23}) lies in $(0, \pi)$, the mirror configuration (robot 3 clockwise of the edge) maps to the same SAS parameters.

The local centroid is:
\begin{equation}\label{eq:appB:28}
x_c^{\text{local}} = \frac{x_1^{\text{local}} + x_2^{\text{local}} + x_3^{\text{local}}}{3}, \quad y_c^{\text{local}} = \frac{y_1^{\text{local}} + y_2^{\text{local}} + y_3^{\text{local}}}{3}
\end{equation}

After centering the triangle at the local origin and applying rotation $\theta_c$ and translation to $(x_c, y_c)$, the global positions become:
\begin{equation}\label{eq:appB:29}
\begin{bmatrix} x_i \\ y_i \end{bmatrix} = \begin{bmatrix} \cos\theta_c & -\sin\theta_c \\ \sin\theta_c & \cos\theta_c \end{bmatrix} \begin{bmatrix} x_i^{\text{local}} - x_c^{\text{local}} \\ y_i^{\text{local}} - y_c^{\text{local}} \end{bmatrix} + \begin{bmatrix} x_c \\ y_c \end{bmatrix}
\end{equation}
for $i \in \{1, 2, 3\}$.

\section{Inverse Jacobian}

The inverse Jacobian maps shape parameter velocities to robot velocities:
\begin{equation}\label{eq:appB:30}
\begin{bmatrix}
\delta x_1 \\ \delta y_1 \\ \delta x_2 \\ \delta y_2 \\ \delta x_3 \\ \delta y_3
\end{bmatrix} =
\mathbf{J}_c^{-1}
\begin{bmatrix}
\delta x_c \\ \delta y_c \\ \delta \theta_c \\ \delta p \\ \delta \beta \\ \delta q
\end{bmatrix}
\end{equation}

where $\mathbf{J}_c^{-1}$ is a $6 \times 6$ matrix containing partial derivatives $\frac{\partial x_i}{\partial \xi_j}$ and $\frac{\partial y_i}{\partial \xi_j}$ for robot positions and shape parameters $\xi_j \in \{x_c, y_c, \theta_c, p, \beta, q\}$. This matrix is computed analytically by taking partial derivatives of the inverse kinematics equations (\ref{eq:appB:29}) with respect to each shape parameter. The explicit entries are not listed here. The simulation software computes them by forward finite differences of the inverse kinematics (\texttt{compute\_inverse\_jacobian} in \texttt{src/control/kinematics.py}), and the \texttt{cluster\_builder} tool generates them in closed symbolic form (\texttt{clusterbuilder.py} with the \texttt{--symbolic} option).
